%%%%%%%%%%%%%%%%%%%%%%%%%%%%%%%%%%%%%%%%%%%%%%%%%%%%%%%%%%%%
% 5.15 POTENTIAL THEORY
% The Composition of Advanced Mathematics
%%%%%%%%%%%%%%%%%%%%%%%%%%%%%%%%%%%%%%%%%%%%%%%%%%%%%%%%%%%%

\section{Potential Theory}
\label{sec:potential-theory}

\prerequisite{
Multivariable calculus, vector calculus, real analysis, measure theory,
Lebesgue integration, complex analysis, functional analysis, calculus
of variations, and distribution theory.
}

\learningobjectives{
By the end of this section, the reader should be able to:
\begin{itemize}
    \item explain the role of harmonic functions in potential theory;
    \item work with the Laplace and Poisson equations;
    \item recognize subharmonic and superharmonic functions;
    \item use the mean-value property;
    \item apply maximum and minimum principles;
    \item understand uniqueness for the Dirichlet problem;
    \item derive and use fundamental solutions of the Laplacian;
    \item construct Newtonian and logarithmic potentials;
    \item distinguish volume, single-layer, and double-layer potentials;
    \item understand Green's identities and Green's functions;
    \item use Poisson kernels for balls and half-spaces;
    \item understand the Dirichlet and Neumann boundary-value problems;
    \item interpret harmonic functions through energy minimization;
    \item understand harmonic measure and the probabilistic viewpoint;
    \item recognize removable singularities of harmonic functions;
    \item understand the connection between harmonic and holomorphic functions;
    \item recognize capacity and polar sets;
    \item connect potential theory with electrostatics, gravitation,
          diffusion, PDEs, probability, and geometric analysis.
\end{itemize}
}

%%%%%%%%%%%%%%%%%%%%%%%%%%%%%%%%%%%%%%%%%%%%%%%%%%%%%%%%%%%%
\subsection{What Is Potential Theory?}
\label{subsec:what-is-potential-theory}
%%%%%%%%%%%%%%%%%%%%%%%%%%%%%%%%%%%%%%%%%%%%%%%%%%%%%%%%%%%%

Potential theory studies functions generated by sources and the
differential equations governing those functions.

The central operator is the Laplacian

\[
\boxed{
\Delta
=
\frac{\partial^2}{\partial x_1^2}
+\cdots+
\frac{\partial^2}{\partial x_n^2}.
}
\]

Potential theory is built primarily around equations such as

\[
\boxed{
\Delta u=0
}
\]

and

\[
\boxed{
-\Delta u=f.
}
\]

These are the Laplace and Poisson equations.

The same mathematical structure appears in

\[
\boxed{
\text{electrostatics},
\quad
\text{gravitation},
\quad
\text{steady-state heat flow},
\quad
\text{fluid mechanics},
\quad
\text{probability}.
}
\]

%%%%%%%%%%%%%%%%%%%%%%%%%%%%%%%%%%%%%%%%%%%%%%%%%%%%%%%%%%%%
\subsection{The Laplacian}
\label{subsec:potential-laplacian}
%%%%%%%%%%%%%%%%%%%%%%%%%%%%%%%%%%%%%%%%%%%%%%%%%%%%%%%%%%%%

For a twice differentiable function

\[
u:\Omega\subseteq\R^n\to\R,
\]

the Laplacian is

\[
\boxed{
\Delta u
=
\sum_{j=1}^{n}
\frac{\partial^2u}{\partial x_j^2}.
}
\]

In three dimensions,

\[
\boxed{
\Delta u
=
u_{xx}+u_{yy}+u_{zz}.
}
\]

The Laplacian measures a form of local deviation from average behavior.

This interpretation becomes precise through the mean-value property.

%%%%%%%%%%%%%%%%%%%%%%%%%%%%%%%%%%%%%%%%%%%%%%%%%%%%%%%%%%%%
\subsection{Harmonic Functions}
\label{subsec:harmonic-functions-potential}
%%%%%%%%%%%%%%%%%%%%%%%%%%%%%%%%%%%%%%%%%%%%%%%%%%%%%%%%%%%%

\begin{definitionbox}{Harmonic Function}
A twice continuously differentiable function

\[
u:\Omega\to\R
\]

is harmonic on an open set \(\Omega\subseteq\R^n\) if

\[
\boxed{
\Delta u=0
}
\]

throughout \(\Omega\).
\end{definitionbox}

Thus harmonic functions are solutions of Laplace's equation.

Examples include

\[
\boxed{
u(x,y)=x,
}
\]

\[
\boxed{
u(x,y)=x^2-y^2,
}
\]

and

\[
\boxed{
u(x,y)=xy.
}
\]

%%%%%%%%%%%%%%%%%%%%%%%%%%%%%%%%%%%%%%%%%%%%%%%%%%%%%%%%%%%%
\subsection{Testing Harmonicity}
\label{subsec:testing-harmonicity}
%%%%%%%%%%%%%%%%%%%%%%%%%%%%%%%%%%%%%%%%%%%%%%%%%%%%%%%%%%%%

Consider

\[
u(x,y)=x^2-y^2.
\]

Then

\[
u_{xx}=2
\]

and

\[
u_{yy}=-2.
\]

Therefore,

\[
\Delta u
=
2-2
=
0.
\]

Hence

\[
\boxed{
u(x,y)=x^2-y^2
}
\]

is harmonic on \(\R^2\).

%%%%%%%%%%%%%%%%%%%%%%%%%%%%%%%%%%%%%%%%%%%%%%%%%%%%%%%%%%%%
\subsection{Poisson's Equation}
\label{subsec:potential-poisson-equation}
%%%%%%%%%%%%%%%%%%%%%%%%%%%%%%%%%%%%%%%%%%%%%%%%%%%%%%%%%%%%

The nonhomogeneous version of Laplace's equation is Poisson's equation.

A common sign convention is

\[
\boxed{
-\Delta u=f.
}
\]

Here \(f\) represents a source density.

For example,

\[
\boxed{
-\Delta u=\rho
}
\]

may describe a potential generated by a mass or charge density
\(\rho\).

Sign conventions vary between fields, so the convention must always be
stated.

%%%%%%%%%%%%%%%%%%%%%%%%%%%%%%%%%%%%%%%%%%%%%%%%%%%%%%%%%%%%
\subsection{Mean-Value Property}
\label{subsec:harmonic-mean-value-property}
%%%%%%%%%%%%%%%%%%%%%%%%%%%%%%%%%%%%%%%%%%%%%%%%%%%%%%%%%%%%

One of the defining features of harmonic functions is that their value
at the center of a ball equals their average over the ball or its
boundary.

\begin{theorem}[Mean-Value Property]
If \(u\) is harmonic on an open set containing the closed ball
\(\overline{B_r(x_0)}\), then

\[
\boxed{
u(x_0)
=
\frac{1}{|\partial B_r|}
\int_{\partial B_r(x_0)}
u\,dS.
}
\]

Also,

\[
\boxed{
u(x_0)
=
\frac{1}{|B_r|}
\int_{B_r(x_0)}
u(x)\,dx.
}
\]
\end{theorem}

Thus a harmonic function has no isolated local spikes or pits.

%%%%%%%%%%%%%%%%%%%%%%%%%%%%%%%%%%%%%%%%%%%%%%%%%%%%%%%%%%%%
\subsection{Harmonic Functions Are Determined by Averages}
\label{subsec:harmonic-average-interpretation}
%%%%%%%%%%%%%%%%%%%%%%%%%%%%%%%%%%%%%%%%%%%%%%%%%%%%%%%%%%%%

The mean-value property gives the intuitive principle

\[
\boxed{
\text{value at a point}
=
\text{average of nearby values}.
}
\]

Consequently, harmonic functions are highly constrained.

Local behavior is controlled by surrounding values, which leads to
powerful maximum principles and regularity results.

%%%%%%%%%%%%%%%%%%%%%%%%%%%%%%%%%%%%%%%%%%%%%%%%%%%%%%%%%%%%
\subsection{Converse Mean-Value Theorem}
\label{subsec:converse-mean-value}
%%%%%%%%%%%%%%%%%%%%%%%%%%%%%%%%%%%%%%%%%%%%%%%%%%%%%%%%%%%%

Under suitable regularity assumptions, the mean-value property also
characterizes harmonic functions.

\begin{theorem}[Converse Mean-Value Principle]
If a continuous function satisfies the mean-value property on every
sufficiently small ball contained in \(\Omega\), then it is harmonic in
\(\Omega\).
\end{theorem}

Thus

\[
\boxed{
\text{harmonicity}
\Longleftrightarrow
\text{local mean-value behavior}
}
\]

under appropriate hypotheses.

%%%%%%%%%%%%%%%%%%%%%%%%%%%%%%%%%%%%%%%%%%%%%%%%%%%%%%%%%%%%
\subsection{Maximum Principle}
\label{subsec:harmonic-maximum-principle}
%%%%%%%%%%%%%%%%%%%%%%%%%%%%%%%%%%%%%%%%%%%%%%%%%%%%%%%%%%%%

\begin{theorem}[Strong Maximum Principle]
Let \(\Omega\) be connected and let \(u\) be harmonic on \(\Omega\).

If \(u\) attains a maximum at an interior point, then \(u\) is constant.
\end{theorem}

Similarly, if a harmonic function attains an interior minimum, it must be
constant.

Therefore a nonconstant harmonic function cannot have a strict interior
maximum or minimum.

%%%%%%%%%%%%%%%%%%%%%%%%%%%%%%%%%%%%%%%%%%%%%%%%%%%%%%%%%%%%
\subsection{Boundary Maximum Principle}
\label{subsec:boundary-maximum-principle}
%%%%%%%%%%%%%%%%%%%%%%%%%%%%%%%%%%%%%%%%%%%%%%%%%%%%%%%%%%%%

Suppose \(\Omega\) is bounded and \(u\) is harmonic in \(\Omega\) and
continuous on \(\overline{\Omega}\).

Then

\[
\boxed{
\max_{\overline{\Omega}}u
=
\max_{\partial\Omega}u
}
\]

and

\[
\boxed{
\min_{\overline{\Omega}}u
=
\min_{\partial\Omega}u.
}
\]

Thus the extreme values are controlled entirely by the boundary.

%%%%%%%%%%%%%%%%%%%%%%%%%%%%%%%%%%%%%%%%%%%%%%%%%%%%%%%%%%%%
\subsection{Uniqueness for the Dirichlet Problem}
\label{subsec:dirichlet-uniqueness}
%%%%%%%%%%%%%%%%%%%%%%%%%%%%%%%%%%%%%%%%%%%%%%%%%%%%%%%%%%%%

Suppose \(u\) and \(v\) are harmonic in \(\Omega\) and satisfy

\[
u=v
\]

on \(\partial\Omega\).

Then

\[
w=u-v
\]

is harmonic and satisfies

\[
w=0
\]

on the boundary.

The maximum principle gives

\[
w\leq0
\]

and

\[
-w\leq0.
\]

Hence

\[
w=0.
\]

Therefore,

\[
\boxed{
u=v.
}
\]

This proves uniqueness for the classical Dirichlet problem whenever a
solution exists.

%%%%%%%%%%%%%%%%%%%%%%%%%%%%%%%%%%%%%%%%%%%%%%%%%%%%%%%%%%%%
\subsection{Dirichlet Problem}
\label{subsec:dirichlet-problem}
%%%%%%%%%%%%%%%%%%%%%%%%%%%%%%%%%%%%%%%%%%%%%%%%%%%%%%%%%%%%

The Dirichlet problem asks for a harmonic function with prescribed
boundary values:

\[
\boxed{
\begin{cases}
\Delta u=0,
& x\in\Omega,\\
u=g,
& x\in\partial\Omega.
\end{cases}
}
\]

The main questions are

\[
\boxed{
\text{existence},
\quad
\text{uniqueness},
\quad
\text{regularity}.
}
\]

Uniqueness follows from the maximum principle.

Existence depends more delicately on the domain and boundary data.

%%%%%%%%%%%%%%%%%%%%%%%%%%%%%%%%%%%%%%%%%%%%%%%%%%%%%%%%%%%%
\subsection{Neumann Problem}
\label{subsec:neumann-problem}
%%%%%%%%%%%%%%%%%%%%%%%%%%%%%%%%%%%%%%%%%%%%%%%%%%%%%%%%%%%%

The Neumann problem prescribes the normal derivative:

\[
\boxed{
\begin{cases}
\Delta u=f,
& x\in\Omega,\\
\dfrac{\partial u}{\partial n}=g,
& x\in\partial\Omega.
\end{cases}
}
\]

Here

\[
\frac{\partial u}{\partial n}
=
\nabla u\cdot n,
\]

where \(n\) is the outward unit normal.

For the Poisson equation, integrating over \(\Omega\) and using the
divergence theorem gives the compatibility condition

\[
\boxed{
\int_\Omega f\,dx
=
\int_{\partial\Omega}g\,dS.
}
\]

For the homogeneous Laplace equation this becomes

\[
\boxed{
\int_{\partial\Omega}g\,dS=0.
}
\]

Solutions of the Neumann problem are generally unique only up to an
additive constant.

%%%%%%%%%%%%%%%%%%%%%%%%%%%%%%%%%%%%%%%%%%%%%%%%%%%%%%%%%%%%
\subsection{Green's First Identity}
\label{subsec:greens-first-identity-potential}
%%%%%%%%%%%%%%%%%%%%%%%%%%%%%%%%%%%%%%%%%%%%%%%%%%%%%%%%%%%%

For sufficiently smooth \(u\) and \(v\),

\[
\boxed{
\int_\Omega
\left(
\nabla u\cdot\nabla v
+
u\Delta v
\right)\,dx
=
\int_{\partial\Omega}
u\frac{\partial v}{\partial n}\,dS.
}
\]

This follows from the divergence identity

\[
\nabla\cdot(u\nabla v)
=
\nabla u\cdot\nabla v
+
u\Delta v.
\]

Green's first identity is fundamental in potential theory and weak PDE
formulations.

%%%%%%%%%%%%%%%%%%%%%%%%%%%%%%%%%%%%%%%%%%%%%%%%%%%%%%%%%%%%
\subsection{Green's Second Identity}
\label{subsec:greens-second-identity-potential}
%%%%%%%%%%%%%%%%%%%%%%%%%%%%%%%%%%%%%%%%%%%%%%%%%%%%%%%%%%%%

Subtracting Green's first identity with \(u\) and \(v\) exchanged gives

\[
\boxed{
\int_\Omega
\left(
u\Delta v-v\Delta u
\right)\,dx
=
\int_{\partial\Omega}
\left(
u\frac{\partial v}{\partial n}
-
v\frac{\partial u}{\partial n}
\right)dS.
}
\]

Green's second identity is the higher-dimensional analogue of integration
by parts used to derive representation formulas.

%%%%%%%%%%%%%%%%%%%%%%%%%%%%%%%%%%%%%%%%%%%%%%%%%%%%%%%%%%%%
\subsection{Fundamental Solution of the Laplacian}
\label{subsec:potential-fundamental-solution}
%%%%%%%%%%%%%%%%%%%%%%%%%%%%%%%%%%%%%%%%%%%%%%%%%%%%%%%%%%%%

A fundamental solution \(\Phi\) satisfies

\[
\boxed{
\Delta\Phi=\delta_0
}
\]

in the distributional sense.

Let

\[
\omega_n
=
|\mathbb{S}^{n-1}|
\]

denote the surface area of the unit sphere in \(\R^n\).

For

\[
n\geq3,
\]

one may use

\[
\boxed{
\Phi(x)
=
\frac{1}{(2-n)\omega_n}
|x|^{2-n}.
}
\]

For

\[
n=2,
\]

one may use

\[
\boxed{
\Phi(x)
=
\frac{1}{2\pi}\log|x|.
}
\]

These formulas correspond to the convention

\[
\Delta\Phi=\delta_0.
\]

For the operator \(-\Delta\), the signs are reversed.

%%%%%%%%%%%%%%%%%%%%%%%%%%%%%%%%%%%%%%%%%%%%%%%%%%%%%%%%%%%%
\subsection{Radial Harmonic Functions}
\label{subsec:radial-harmonic-functions}
%%%%%%%%%%%%%%%%%%%%%%%%%%%%%%%%%%%%%%%%%%%%%%%%%%%%%%%%%%%%

Suppose

\[
u(x)=v(r),
\qquad
r=|x|.
\]

Then

\[
\boxed{
\Delta u
=
v''(r)
+
\frac{n-1}{r}v'(r).
}
\]

Thus radial harmonic functions satisfy

\[
v''
+
\frac{n-1}{r}v'
=
0.
\]

Equivalently,

\[
\frac{d}{dr}
\left(
r^{n-1}v'(r)
\right)
=
0.
\]

Hence

\[
v'(r)
=
Cr^{1-n}.
\]

For \(n\geq3\),

\[
\boxed{
v(r)
=
A+Br^{2-n}.
}
\]

For \(n=2\),

\[
\boxed{
v(r)
=
A+B\log r.
}
\]

These are precisely the forms appearing in the fundamental solutions.

%%%%%%%%%%%%%%%%%%%%%%%%%%%%%%%%%%%%%%%%%%%%%%%%%%%%%%%%%%%%
\subsection{Newtonian Potential}
\label{subsec:newtonian-potential}
%%%%%%%%%%%%%%%%%%%%%%%%%%%%%%%%%%%%%%%%%%%%%%%%%%%%%%%%%%%%

Suppose

\[
n\geq3
\]

and \(f\) is sufficiently well behaved.

Using the fundamental solution for \(-\Delta\),

\[
\Gamma(x)
=
\frac{1}{(n-2)\omega_n}|x|^{2-n},
\]

define the Newtonian potential

\[
\boxed{
u(x)
=
(\Gamma*f)(x)
=
\int_{\R^n}
\Gamma(x-y)f(y)\,dy.
}
\]

Formally,

\[
-\Delta u
=
(-\Delta\Gamma)*f
=
\delta*f
=
f.
\]

Therefore,

\[
\boxed{
-\Delta u=f.
}
\]

%%%%%%%%%%%%%%%%%%%%%%%%%%%%%%%%%%%%%%%%%%%%%%%%%%%%%%%%%%%%
\subsection{Logarithmic Potential}
\label{subsec:logarithmic-potential}
%%%%%%%%%%%%%%%%%%%%%%%%%%%%%%%%%%%%%%%%%%%%%%%%%%%%%%%%%%%%

In two dimensions, the fundamental solution has logarithmic form.

For the operator \(-\Delta\), one may use

\[
\boxed{
\Gamma(x)
=
-\frac{1}{2\pi}\log|x|.
}
\]

The corresponding potential is

\[
\boxed{
u(x)
=
\int_{\R^2}
\Gamma(x-y)f(y)\,dy.
}
\]

This is called a logarithmic potential.

%%%%%%%%%%%%%%%%%%%%%%%%%%%%%%%%%%%%%%%%%%%%%%%%%%%%%%%%%%%%
\subsection{Electrostatic Interpretation}
\label{subsec:electrostatic-potential}
%%%%%%%%%%%%%%%%%%%%%%%%%%%%%%%%%%%%%%%%%%%%%%%%%%%%%%%%%%%%

In electrostatics, a charge density \(\rho\) generates an electric
potential \(V\).

In suitable units,

\[
\boxed{
-\Delta V=\rho.
}
\]

The electric field is

\[
\boxed{
E=-\nabla V.
}
\]

Thus

\[
\boxed{
\text{charge density}
\longrightarrow
\text{potential}
\longrightarrow
\text{field}.
}
\]

Outside the support of the charge,

\[
\rho=0,
\]

so

\[
\boxed{
\Delta V=0.
}
\]

The potential is harmonic away from its sources.

%%%%%%%%%%%%%%%%%%%%%%%%%%%%%%%%%%%%%%%%%%%%%%%%%%%%%%%%%%%%
\subsection{Gravitational Interpretation}
\label{subsec:gravitational-potential}
%%%%%%%%%%%%%%%%%%%%%%%%%%%%%%%%%%%%%%%%%%%%%%%%%%%%%%%%%%%%

Newtonian gravitation has an analogous structure.

A mass density \(\rho\) generates a gravitational potential.

The gravitational field is obtained from its gradient, with the physical
sign determined by convention.

Thus potential theory unifies gravitational and electrostatic field
models mathematically.

%%%%%%%%%%%%%%%%%%%%%%%%%%%%%%%%%%%%%%%%%%%%%%%%%%%%%%%%%%%%
\subsection{Green's Functions}
\label{subsec:potential-greens-functions}
%%%%%%%%%%%%%%%%%%%%%%%%%%%%%%%%%%%%%%%%%%%%%%%%%%%%%%%%%%%%

\begin{definitionbox}{Green's Function}
A Green's function for a differential operator \(L\) on a domain
\(\Omega\) is a kernel \(G(x,y)\) satisfying, in the appropriate sense,

\[
\boxed{
L_xG(x,y)=\delta_y(x),
}
\]

together with specified boundary conditions.
\end{definitionbox}

For the Dirichlet Laplacian, one commonly chooses

\[
\boxed{
G(x,y)=0
\qquad
\text{for }x\in\partial\Omega.
}
\]

Green's functions combine

\[
\boxed{
\text{fundamental solution}
+
\text{boundary correction}.
}
\]

%%%%%%%%%%%%%%%%%%%%%%%%%%%%%%%%%%%%%%%%%%%%%%%%%%%%%%%%%%%%
\subsection{Green Representation Formula}
\label{subsec:green-representation-formula}
%%%%%%%%%%%%%%%%%%%%%%%%%%%%%%%%%%%%%%%%%%%%%%%%%%%%%%%%%%%%

Suppose \(G\) is a Dirichlet Green's function satisfying

\[
-\Delta_xG(x,y)=\delta_y(x)
\]

and

\[
G(x,y)=0
\]

for \(x\in\partial\Omega\).

For a sufficiently smooth solution of

\[
-\Delta u=f,
\]

Green's identity yields a representation of the form

\[
\boxed{
u(y)
=
\int_\Omega
G(x,y)f(x)\,dx
-
\int_{\partial\Omega}
u(x)
\frac{\partial G}{\partial n_x}(x,y)\,dS_x.
}
\]

The precise signs depend on the chosen convention for the Laplacian and
Green's function.

Representation formulas convert differential equations into integral
formulas.

%%%%%%%%%%%%%%%%%%%%%%%%%%%%%%%%%%%%%%%%%%%%%%%%%%%%%%%%%%%%
\subsection{Poisson Kernel for the Unit Disk}
\label{subsec:poisson-kernel-disk}
%%%%%%%%%%%%%%%%%%%%%%%%%%%%%%%%%%%%%%%%%%%%%%%%%%%%%%%%%%%%

For the unit disk

\[
\mathbb{D}
=
\{z\in\mathbb{C}:|z|<1\},
\]

the Poisson kernel is

\[
\boxed{
P_r(\theta)
=
\frac{1-r^2}
{1-2r\cos\theta+r^2},
\qquad
0\leq r<1.
}
\]

If boundary data are given by \(g\), then the harmonic extension is

\[
\boxed{
u(re^{i\theta})
=
\frac{1}{2\pi}
\int_0^{2\pi}
P_r(\theta-t)g(e^{it})\,dt.
}
\]

This is the Poisson integral formula.

%%%%%%%%%%%%%%%%%%%%%%%%%%%%%%%%%%%%%%%%%%%%%%%%%%%%%%%%%%%%
\subsection{Properties of the Poisson Kernel}
\label{subsec:poisson-kernel-properties}
%%%%%%%%%%%%%%%%%%%%%%%%%%%%%%%%%%%%%%%%%%%%%%%%%%%%%%%%%%%%

The Poisson kernel satisfies

\[
\boxed{
P_r(\theta)>0
}
\]

and

\[
\boxed{
\frac{1}{2\pi}
\int_0^{2\pi}
P_r(\theta)\,d\theta
=
1.
}
\]

As

\[
r\to1^-,
\]

the kernel concentrates near

\[
\theta=0.
\]

Thus the Poisson kernels form an approximate identity on the circle.

This connects potential theory with Fourier and harmonic analysis.

%%%%%%%%%%%%%%%%%%%%%%%%%%%%%%%%%%%%%%%%%%%%%%%%%%%%%%%%%%%%
\subsection{Poisson Kernel for the Upper Half-Space}
\label{subsec:poisson-kernel-half-space}
%%%%%%%%%%%%%%%%%%%%%%%%%%%%%%%%%%%%%%%%%%%%%%%%%%%%%%%%%%%%

Let

\[
\R_+^n
=
\{
(x',t)\in\R^{n-1}\times\R:t>0
\}.
\]

The Poisson kernel for the upper half-space has the form

\[
\boxed{
P_t(x')
=
c_n
\frac{t}
{\left(|x'|^2+t^2\right)^{n/2}},
}
\]

where \(c_n\) is chosen so that

\[
\boxed{
\int_{\R^{n-1}}P_t(x')\,dx'=1.
}
\]

Given suitable boundary data \(f\),

\[
\boxed{
u(x',t)
=
(P_t*f)(x')
}
\]

is harmonic in the upper half-space.

%%%%%%%%%%%%%%%%%%%%%%%%%%%%%%%%%%%%%%%%%%%%%%%%%%%%%%%%%%%%
\subsection{Poisson Kernel for a Ball}
\label{subsec:poisson-kernel-ball}
%%%%%%%%%%%%%%%%%%%%%%%%%%%%%%%%%%%%%%%%%%%%%%%%%%%%%%%%%%%%

For the ball

\[
B_R(0)\subseteq\R^n,
\]

the Poisson kernel can be written

\[
\boxed{
P(x,\xi)
=
\frac{R^2-|x|^2}
{\omega_nR|x-\xi|^n},
\qquad
|x|<R,
\quad
|\xi|=R.
}
\]

If \(g\) is suitable boundary data, then

\[
\boxed{
u(x)
=
\int_{\partial B_R}
P(x,\xi)g(\xi)\,dS_\xi
}
\]

is harmonic inside the ball and assumes the prescribed boundary values
in the appropriate sense.

%%%%%%%%%%%%%%%%%%%%%%%%%%%%%%%%%%%%%%%%%%%%%%%%%%%%%%%%%%%%
\subsection{Worked Example: Harmonic Extension on the Disk}
\label{subsec:worked-harmonic-extension-disk}
%%%%%%%%%%%%%%%%%%%%%%%%%%%%%%%%%%%%%%%%%%%%%%%%%%%%%%%%%%%%

\begin{workedexamplebox}{Boundary Data \(g(\theta)=\cos\theta\)}

Consider the unit disk with boundary values

\[
g(\theta)=\cos\theta.
\]

The harmonic function

\[
u(r,\theta)=r\cos\theta
\]

has boundary values

\[
u(1,\theta)=\cos\theta.
\]

Since

\[
r\cos\theta=x,
\]

the function is simply

\[
u(x,y)=x.
\]

Because

\[
\Delta x=0,
\]

it is harmonic.

\begin{answerbox}
\[
\boxed{
u(r,\theta)=r\cos\theta.
}
\]
\end{answerbox}

\end{workedexamplebox}

%%%%%%%%%%%%%%%%%%%%%%%%%%%%%%%%%%%%%%%%%%%%%%%%%%%%%%%%%%%%
\subsection{Harnack Inequality}
\label{subsec:harnack-inequality}
%%%%%%%%%%%%%%%%%%%%%%%%%%%%%%%%%%%%%%%%%%%%%%%%%%%%%%%%%%%%

Positive harmonic functions satisfy strong comparison estimates.

\begin{theorem}[Harnack Inequality --- Qualitative Form]
Let \(u\) be a positive harmonic function on a domain containing a compact
set \(K\).

Then values of \(u\) on \(K\) are comparable:

\[
\boxed{
\sup_Ku
\leq
C\inf_Ku,
}
\]

where \(C\) depends on the geometry of \(K\) and its position inside the
domain.
\end{theorem}

This means a positive harmonic function cannot vary arbitrarily rapidly
away from the boundary.

%%%%%%%%%%%%%%%%%%%%%%%%%%%%%%%%%%%%%%%%%%%%%%%%%%%%%%%%%%%%
\subsection{Harnack Convergence Principle}
\label{subsec:harnack-convergence}
%%%%%%%%%%%%%%%%%%%%%%%%%%%%%%%%%%%%%%%%%%%%%%%%%%%%%%%%%%%%

Sequences of harmonic functions exhibit strong compactness behavior.

A typical Harnack convergence principle states that a locally bounded
monotone sequence of harmonic functions either converges harmonically or
diverges in a controlled manner.

Such results are early examples of compactness principles in elliptic
PDE theory.

%%%%%%%%%%%%%%%%%%%%%%%%%%%%%%%%%%%%%%%%%%%%%%%%%%%%%%%%%%%%
\subsection{Smoothness of Harmonic Functions}
\label{subsec:harmonic-smoothness}
%%%%%%%%%%%%%%%%%%%%%%%%%%%%%%%%%%%%%%%%%%%%%%%%%%%%%%%%%%%%

The equation

\[
\Delta u=0
\]

appears to involve only second derivatives.

Nevertheless, harmonic functions are much smoother.

\begin{theorem}[Regularity of Harmonic Functions]
Every harmonic function is infinitely differentiable.

Indeed, harmonic functions are real analytic.
\end{theorem}

Thus

\[
\boxed{
\Delta u=0
\Longrightarrow
u\in C^\infty
}
\]

and much more.

This phenomenon is called elliptic regularity.

%%%%%%%%%%%%%%%%%%%%%%%%%%%%%%%%%%%%%%%%%%%%%%%%%%%%%%%%%%%%
\subsection{Derivative Estimates}
\label{subsec:harmonic-derivative-estimates}
%%%%%%%%%%%%%%%%%%%%%%%%%%%%%%%%%%%%%%%%%%%%%%%%%%%%%%%%%%%%

The mean-value property leads to estimates for derivatives.

For a harmonic function on a ball, derivatives at an interior point can
be bounded in terms of the size of the function on a larger surrounding
ball.

Schematically,

\[
\boxed{
|D^\alpha u(x_0)|
\leq
\frac{C_\alpha}{r^{|\alpha|}}
\sup_{B_r(x_0)}|u|.
}
\]

These estimates help prove analyticity and compactness properties.

%%%%%%%%%%%%%%%%%%%%%%%%%%%%%%%%%%%%%%%%%%%%%%%%%%%%%%%%%%%%
\subsection{Liouville Theorem for Harmonic Functions}
\label{subsec:harmonic-liouville}
%%%%%%%%%%%%%%%%%%%%%%%%%%%%%%%%%%%%%%%%%%%%%%%%%%%%%%%%%%%%

\begin{theorem}[Liouville Theorem]
Every bounded harmonic function

\[
u:\R^n\to\R
\]

is constant.
\end{theorem}

This parallels Liouville's theorem in complex analysis.

The result follows from interior estimates or the mean-value property.

%%%%%%%%%%%%%%%%%%%%%%%%%%%%%%%%%%%%%%%%%%%%%%%%%%%%%%%%%%%%
\subsection{Harmonic and Holomorphic Functions}
\label{subsec:harmonic-holomorphic}
%%%%%%%%%%%%%%%%%%%%%%%%%%%%%%%%%%%%%%%%%%%%%%%%%%%%%%%%%%%%

Suppose

\[
f(z)=u(x,y)+iv(x,y)
\]

is holomorphic.

The Cauchy--Riemann equations imply

\[
\boxed{
\Delta u=0
}
\]

and

\[
\boxed{
\Delta v=0.
}
\]

Therefore the real and imaginary parts of a holomorphic function are
harmonic.

This gives a major bridge between Complex Analysis and Potential Theory.

%%%%%%%%%%%%%%%%%%%%%%%%%%%%%%%%%%%%%%%%%%%%%%%%%%%%%%%%%%%%
\subsection{Harmonic Conjugates}
\label{subsec:potential-harmonic-conjugates}
%%%%%%%%%%%%%%%%%%%%%%%%%%%%%%%%%%%%%%%%%%%%%%%%%%%%%%%%%%%%

If \(u\) is harmonic on a simply connected planar domain, then under the
standard hypotheses there exists a harmonic function \(v\) such that

\[
\boxed{
f=u+iv
}
\]

is holomorphic.

The function \(v\) is called a harmonic conjugate of \(u\).

Thus two-dimensional potential theory is deeply connected with complex
analytic methods.

%%%%%%%%%%%%%%%%%%%%%%%%%%%%%%%%%%%%%%%%%%%%%%%%%%%%%%%%%%%%
\subsection{Subharmonic Functions}
\label{subsec:subharmonic-functions}
%%%%%%%%%%%%%%%%%%%%%%%%%%%%%%%%%%%%%%%%%%%%%%%%%%%%%%%%%%%%

\begin{definitionbox}{Subharmonic Function}
A sufficiently regular function \(u\) is subharmonic if

\[
\boxed{
\Delta u\geq0.
}
\]
\end{definitionbox}

For less regular functions, subharmonicity is defined using appropriate
mean-value inequalities or distributional conditions.

A subharmonic function satisfies the heuristic inequality

\[
\boxed{
u(x_0)
\leq
\text{average of \(u\) around \(x_0\)}.
}
\]

%%%%%%%%%%%%%%%%%%%%%%%%%%%%%%%%%%%%%%%%%%%%%%%%%%%%%%%%%%%%
\subsection{Superharmonic Functions}
\label{subsec:superharmonic-functions}
%%%%%%%%%%%%%%%%%%%%%%%%%%%%%%%%%%%%%%%%%%%%%%%%%%%%%%%%%%%%

\begin{definitionbox}{Superharmonic Function}
A sufficiently regular function \(u\) is superharmonic if

\[
\boxed{
\Delta u\leq0.
}
\]
\end{definitionbox}

Equivalently,

\[
\boxed{
u\text{ is superharmonic}
\Longleftrightarrow
-u\text{ is subharmonic}.
}
\]

Harmonic functions are simultaneously subharmonic and superharmonic.

%%%%%%%%%%%%%%%%%%%%%%%%%%%%%%%%%%%%%%%%%%%%%%%%%%%%%%%%%%%%
\subsection{Comparison Principle}
\label{subsec:potential-comparison-principle}
%%%%%%%%%%%%%%%%%%%%%%%%%%%%%%%%%%%%%%%%%%%%%%%%%%%%%%%%%%%%

Suppose

\[
\Delta u\geq\Delta v
\]

in \(\Omega\) and

\[
u\leq v
\]

on \(\partial\Omega\), with sign conventions chosen consistently.

Applying the maximum principle to the difference produces comparison
results.

Such principles allow solutions to be bounded above and below by simpler
comparison functions.

%%%%%%%%%%%%%%%%%%%%%%%%%%%%%%%%%%%%%%%%%%%%%%%%%%%%%%%%%%%%
\subsection{Barrier Functions}
\label{subsec:barrier-functions}
%%%%%%%%%%%%%%%%%%%%%%%%%%%%%%%%%%%%%%%%%%%%%%%%%%%%%%%%%%%%

A barrier is a specially constructed subharmonic or superharmonic
function used to control a solution near a boundary point.

Barriers are useful for proving that solutions of the Dirichlet problem
actually approach their prescribed boundary values.

Boundary geometry therefore plays a central role in existence and
regularity theory.

%%%%%%%%%%%%%%%%%%%%%%%%%%%%%%%%%%%%%%%%%%%%%%%%%%%%%%%%%%%%
\subsection{Perron's Method}
\label{subsec:perrons-method}
%%%%%%%%%%%%%%%%%%%%%%%%%%%%%%%%%%%%%%%%%%%%%%%%%%%%%%%%%%%%

Perron's method constructs harmonic functions using families of
subharmonic or superharmonic functions.

Rather than explicitly solving Laplace's equation, one forms an envelope
of comparison functions satisfying appropriate boundary restrictions.

Schematically,

\[
\boxed{
u(x)
=
\sup\{
v(x):v\text{ belongs to an admissible subharmonic family}
\}.
}
\]

Under suitable conditions, the resulting function solves the Dirichlet
problem.

%%%%%%%%%%%%%%%%%%%%%%%%%%%%%%%%%%%%%%%%%%%%%%%%%%%%%%%%%%%%
\subsection{Removable Singularities}
\label{subsec:harmonic-removable-singularities}
%%%%%%%%%%%%%%%%%%%%%%%%%%%%%%%%%%%%%%%%%%%%%%%%%%%%%%%%%%%%

Suppose \(u\) is harmonic on

\[
\Omega\setminus\{a\}.
\]

Under suitable boundedness conditions near \(a\), the apparent
singularity may be removable.

A typical result states that if \(u\) is harmonic near \(a\), except
possibly at \(a\), and remains bounded near \(a\), then \(u\) extends
harmonically across \(a\).

Thus isolated singularities of harmonic functions are constrained by
their growth.

%%%%%%%%%%%%%%%%%%%%%%%%%%%%%%%%%%%%%%%%%%%%%%%%%%%%%%%%%%%%
\subsection{Classification of Radial Singularities}
\label{subsec:radial-singularities}
%%%%%%%%%%%%%%%%%%%%%%%%%%%%%%%%%%%%%%%%%%%%%%%%%%%%%%%%%%%%

For \(n\geq3\), the basic isolated harmonic singularity behaves like

\[
\boxed{
|x|^{2-n}.
}
\]

For \(n=2\), it behaves like

\[
\boxed{
\log|x|.
}
\]

These are exactly the singularities of the fundamental solution.

They play a role analogous to poles and logarithmic singularities in
complex analysis.

%%%%%%%%%%%%%%%%%%%%%%%%%%%%%%%%%%%%%%%%%%%%%%%%%%%%%%%%%%%%
\subsection{Layer Potentials}
\label{subsec:layer-potentials}
%%%%%%%%%%%%%%%%%%%%%%%%%%%%%%%%%%%%%%%%%%%%%%%%%%%%%%%%%%%%

Potential theory can construct harmonic functions from distributions
supported on surfaces.

Let

\[
\Gamma(x-y)
\]

be a fundamental solution for \(-\Delta\).

A single-layer potential has the form

\[
\boxed{
S\sigma(x)
=
\int_{\partial\Omega}
\Gamma(x-y)\sigma(y)\,dS_y.
}
\]

Here \(\sigma\) is a surface density.

%%%%%%%%%%%%%%%%%%%%%%%%%%%%%%%%%%%%%%%%%%%%%%%%%%%%%%%%%%%%
\subsection{Double-Layer Potentials}
\label{subsec:double-layer-potentials}
%%%%%%%%%%%%%%%%%%%%%%%%%%%%%%%%%%%%%%%%%%%%%%%%%%%%%%%%%%%%

A double-layer potential has the form

\[
\boxed{
D\mu(x)
=
\int_{\partial\Omega}
\frac{\partial\Gamma(x-y)}
{\partial n_y}
\mu(y)\,dS_y.
}
\]

Layer potentials convert boundary-value problems into integral equations
on the boundary.

They are fundamental in both theoretical and numerical potential theory.

%%%%%%%%%%%%%%%%%%%%%%%%%%%%%%%%%%%%%%%%%%%%%%%%%%%%%%%%%%%%
\subsection{Jump Relations}
\label{subsec:layer-potential-jump-relations}
%%%%%%%%%%%%%%%%%%%%%%%%%%%%%%%%%%%%%%%%%%%%%%%%%%%%%%%%%%%%

Layer potentials have characteristic behavior as one approaches the
boundary.

Single-layer potentials are typically continuous across sufficiently
regular boundaries, while their normal derivatives may jump.

Double-layer potentials themselves exhibit a jump across the boundary.

These jump relations lead directly to boundary integral equations.

%%%%%%%%%%%%%%%%%%%%%%%%%%%%%%%%%%%%%%%%%%%%%%%%%%%%%%%%%%%%
\subsection{Dirichlet Energy}
\label{subsec:potential-dirichlet-energy}
%%%%%%%%%%%%%%%%%%%%%%%%%%%%%%%%%%%%%%%%%%%%%%%%%%%%%%%%%%%%

For a sufficiently regular function \(u\), define the Dirichlet energy

\[
\boxed{
E[u]
=
\frac12
\int_\Omega
|\nabla u|^2\,dx.
}
\]

If boundary values are fixed, critical points satisfy

\[
\boxed{
\Delta u=0.
}
\]

Thus harmonic functions are variational objects.

%%%%%%%%%%%%%%%%%%%%%%%%%%%%%%%%%%%%%%%%%%%%%%%%%%%%%%%%%%%%
\subsection{Dirichlet Principle}
\label{subsec:potential-dirichlet-principle}
%%%%%%%%%%%%%%%%%%%%%%%%%%%%%%%%%%%%%%%%%%%%%%%%%%%%%%%%%%%%

Suppose one seeks a function satisfying

\[
u=g
\]

on \(\partial\Omega\).

The Dirichlet principle seeks to minimize

\[
\boxed{
E[u]
=
\frac12
\int_\Omega
|\nabla u|^2\,dx
}
\]

over all admissible functions with the prescribed boundary values.

Under an appropriate functional-analytic framework, the minimizer is the
harmonic solution of the Dirichlet problem.

Thus

\[
\boxed{
\text{Laplace equation}
\Longleftrightarrow
\text{energy minimization}.
}
\]

%%%%%%%%%%%%%%%%%%%%%%%%%%%%%%%%%%%%%%%%%%%%%%%%%%%%%%%%%%%%
\subsection{Energy Uniqueness}
\label{subsec:energy-uniqueness}
%%%%%%%%%%%%%%%%%%%%%%%%%%%%%%%%%%%%%%%%%%%%%%%%%%%%%%%%%%%%

Suppose \(u\) satisfies

\[
\Delta u=0
\]

and

\[
u=0
\]

on \(\partial\Omega\).

Green's first identity gives

\[
\int_\Omega
|\nabla u|^2\,dx
=
0.
\]

Therefore,

\[
\nabla u=0.
\]

Hence \(u\) is constant.

Since its boundary value is zero,

\[
\boxed{
u=0.
}
\]

This gives another proof of uniqueness for the Dirichlet problem.

%%%%%%%%%%%%%%%%%%%%%%%%%%%%%%%%%%%%%%%%%%%%%%%%%%%%%%%%%%%%
\subsection{Harmonic Measure}
\label{subsec:harmonic-measure}
%%%%%%%%%%%%%%%%%%%%%%%%%%%%%%%%%%%%%%%%%%%%%%%%%%%%%%%%%%%%

For a domain \(\Omega\) and interior point \(x\), harmonic measure assigns
weights to portions of the boundary.

It is commonly denoted

\[
\boxed{
\omega^x.
}
\]

For suitable boundary data \(g\), a harmonic solution may be represented
as

\[
\boxed{
u(x)
=
\int_{\partial\Omega}
g(\xi)\,d\omega^x(\xi).
}
\]

Thus harmonic measure describes how strongly each part of the boundary
influences the value at \(x\).

%%%%%%%%%%%%%%%%%%%%%%%%%%%%%%%%%%%%%%%%%%%%%%%%%%%%%%%%%%%%
\subsection{Probabilistic Interpretation}
\label{subsec:potential-probabilistic-interpretation}
%%%%%%%%%%%%%%%%%%%%%%%%%%%%%%%%%%%%%%%%%%%%%%%%%%%%%%%%%%%%

Potential theory has a remarkable connection with Brownian motion.

Let

\[
B_t
\]

denote Brownian motion starting at \(x\), and let

\[
\tau_\Omega
\]

be the first exit time from \(\Omega\).

For suitable boundary data \(g\), the solution of the Dirichlet problem
can be represented probabilistically as

\[
\boxed{
u(x)
=
\mathbb{E}_x
\left[
g(B_{\tau_\Omega})
\right].
}
\]

Thus harmonic measure is the exit distribution of Brownian motion.

%%%%%%%%%%%%%%%%%%%%%%%%%%%%%%%%%%%%%%%%%%%%%%%%%%%%%%%%%%%%
\subsection{Mean Value and Brownian Motion}
\label{subsec:mean-value-brownian-motion}
%%%%%%%%%%%%%%%%%%%%%%%%%%%%%%%%%%%%%%%%%%%%%%%%%%%%%%%%%%%%

The probabilistic interpretation explains the mean-value property.

Brownian motion starting at the center of a ball exits the ball with a
symmetric distribution.

The expected boundary value therefore reproduces the harmonic value at
the center.

This provides a bridge between

\[
\boxed{
\text{PDEs},
\quad
\text{potential theory},
\quad
\text{probability theory}.
}
\]

%%%%%%%%%%%%%%%%%%%%%%%%%%%%%%%%%%%%%%%%%%%%%%%%%%%%%%%%%%%%
\subsection{Capacity}
\label{subsec:potential-capacity}
%%%%%%%%%%%%%%%%%%%%%%%%%%%%%%%%%%%%%%%%%%%%%%%%%%%%%%%%%%%%

Potential theory assigns a notion of size to sets that differs from
length, area, or volume.

This quantity is called capacity.

One variational form of capacity for a compact set \(K\subseteq\R^n\)
is based on minimizing an energy among functions satisfying

\[
u\geq1
\]

near \(K\).

Schematically,

\[
\boxed{
\operatorname{Cap}(K)
=
\inf
\int_{\R^n}
|\nabla u|^2\,dx
}
\]

over an appropriate admissible class.

Exact definitions depend on the dimension and potential-theoretic
setting.

%%%%%%%%%%%%%%%%%%%%%%%%%%%%%%%%%%%%%%%%%%%%%%%%%%%%%%%%%%%%
\subsection{Why Capacity Matters}
\label{subsec:why-capacity-matters}
%%%%%%%%%%%%%%%%%%%%%%%%%%%%%%%%%%%%%%%%%%%%%%%%%%%%%%%%%%%%

A set can have measure zero while still having nontrivial
potential-theoretic influence.

Capacity measures whether a set is large enough to affect harmonic
functions or boundary-value problems.

It is therefore particularly useful for studying

\[
\boxed{
\text{exceptional sets},
\quad
\text{removable singularities},
\quad
\text{fine boundary behavior}.
}
\]

%%%%%%%%%%%%%%%%%%%%%%%%%%%%%%%%%%%%%%%%%%%%%%%%%%%%%%%%%%%%
\subsection{Polar Sets}
\label{subsec:polar-sets}
%%%%%%%%%%%%%%%%%%%%%%%%%%%%%%%%%%%%%%%%%%%%%%%%%%%%%%%%%%%%

Very small sets from the viewpoint of potential theory are called polar
sets.

Informally, a polar set can be contained in the set where a suitable
superharmonic function becomes infinite.

Such sets often have zero capacity.

The concept becomes important in fine potential theory and probabilistic
questions involving Brownian motion.

%%%%%%%%%%%%%%%%%%%%%%%%%%%%%%%%%%%%%%%%%%%%%%%%%%%%%%%%%%%%
\subsection{Equilibrium Measures}
\label{subsec:equilibrium-measures}
%%%%%%%%%%%%%%%%%%%%%%%%%%%%%%%%%%%%%%%%%%%%%%%%%%%%%%%%%%%%

Potential theory can associate a measure to a set by minimizing an
interaction energy.

For example, for a kernel \(K(x,y)\), one may study

\[
\boxed{
I(\mu)
=
\iint
K(x,y)\,d\mu(x)\,d\mu(y).
}
\]

A measure minimizing the energy under a normalization constraint is
called an equilibrium measure.

This connects potential theory with measure theory, approximation theory,
and logarithmic energy problems.

%%%%%%%%%%%%%%%%%%%%%%%%%%%%%%%%%%%%%%%%%%%%%%%%%%%%%%%%%%%%
\subsection{The Method of Images}
\label{subsec:method-of-images}
%%%%%%%%%%%%%%%%%%%%%%%%%%%%%%%%%%%%%%%%%%%%%%%%%%%%%%%%%%%%

For highly symmetric domains, Green's functions can sometimes be
constructed by introducing fictitious reflected sources.

This technique is called the method of images.

For example, a point source near a flat boundary can be paired with an
appropriately signed reflected source so that the desired boundary
condition is satisfied.

The method is widely used in electrostatics.

%%%%%%%%%%%%%%%%%%%%%%%%%%%%%%%%%%%%%%%%%%%%%%%%%%%%%%%%%%%%
\subsection{Separation of Variables}
\label{subsec:potential-separation-variables}
%%%%%%%%%%%%%%%%%%%%%%%%%%%%%%%%%%%%%%%%%%%%%%%%%%%%%%%%%%%%

On symmetric domains, Laplace's equation can often be solved by assuming
a product form.

For example,

\[
u(x,y)=X(x)Y(y).
\]

Substituting into

\[
u_{xx}+u_{yy}=0
\]

gives

\[
\frac{X''}{X}
=
-
\frac{Y''}{Y}.
\]

Since one side depends only on \(x\) and the other only on \(y\), both
must equal a constant.

This produces ordinary differential equations whose solutions can be
combined into harmonic expansions.

%%%%%%%%%%%%%%%%%%%%%%%%%%%%%%%%%%%%%%%%%%%%%%%%%%%%%%%%%%%%
\subsection{Fourier Series and Boundary Data}
\label{subsec:potential-fourier-boundary}
%%%%%%%%%%%%%%%%%%%%%%%%%%%%%%%%%%%%%%%%%%%%%%%%%%%%%%%%%%%%

For rectangular or circular domains, boundary data can often be expanded
into Fourier series.

Each Fourier mode has a corresponding harmonic extension.

Thus the solution can be assembled mode by mode.

This is one of the fundamental historical connections between

\[
\boxed{
\text{Fourier analysis}
}
\]

and

\[
\boxed{
\text{potential theory}.
}
\]

%%%%%%%%%%%%%%%%%%%%%%%%%%%%%%%%%%%%%%%%%%%%%%%%%%%%%%%%%%%%
\subsection{Spherical Harmonics}
\label{subsec:potential-spherical-harmonics}
%%%%%%%%%%%%%%%%%%%%%%%%%%%%%%%%%%%%%%%%%%%%%%%%%%%%%%%%%%%%

In dimensions involving spherical geometry, harmonic functions can be
expanded using spherical harmonics.

A spherical harmonic is the restriction to the sphere of a homogeneous
harmonic polynomial.

They satisfy eigenvalue equations for the Laplace--Beltrami operator on
the sphere.

Spherical harmonics are fundamental in

\[
\boxed{
\text{gravitation},
\quad
\text{electromagnetism},
\quad
\text{quantum mechanics},
\quad
\text{geophysics}.
}
\]

%%%%%%%%%%%%%%%%%%%%%%%%%%%%%%%%%%%%%%%%%%%%%%%%%%%%%%%%%%%%
\subsection{Potential Theory on Manifolds}
\label{subsec:potential-theory-manifolds}
%%%%%%%%%%%%%%%%%%%%%%%%%%%%%%%%%%%%%%%%%%%%%%%%%%%%%%%%%%%%

On a Riemannian manifold, the Euclidean Laplacian is replaced by the
Laplace--Beltrami operator

\[
\boxed{
\Delta_g.
}
\]

One then studies

\[
\boxed{
\Delta_g u=0.
}
\]

This extends harmonic-function theory to curved spaces.

The resulting theory connects Potential Theory with Differential Geometry
and Geometric Analysis.

%%%%%%%%%%%%%%%%%%%%%%%%%%%%%%%%%%%%%%%%%%%%%%%%%%%%%%%%%%%%
\subsection{Potential Theory and Elliptic PDEs}
\label{subsec:potential-elliptic-pdes}
%%%%%%%%%%%%%%%%%%%%%%%%%%%%%%%%%%%%%%%%%%%%%%%%%%%%%%%%%%%%

The Laplacian is the model elliptic differential operator.

Many ideas first developed for harmonic functions extend to equations of
the form

\[
\boxed{
Lu=f,
}
\]

where \(L\) is a second-order elliptic operator.

Important themes include

\[
\boxed{
\text{maximum principles},
\quad
\text{regularity},
\quad
\text{Green's functions},
\quad
\text{boundary behavior}.
}
\]

Potential theory therefore forms a major foundation for elliptic PDE
theory.

%%%%%%%%%%%%%%%%%%%%%%%%%%%%%%%%%%%%%%%%%%%%%%%%%%%%%%%%%%%%
\subsection{Potential Theory and Distribution Theory}
\label{subsec:potential-distribution-connection}
%%%%%%%%%%%%%%%%%%%%%%%%%%%%%%%%%%%%%%%%%%%%%%%%%%%%%%%%%%%%

Fundamental solutions are singular at their source points.

For example,

\[
|x|^{2-n}
\]

is not classically harmonic at \(x=0\).

Distributionally,

\[
\boxed{
-\Delta\Gamma=\delta_0.
}
\]

Thus Distribution Theory provides the correct framework for interpreting
point sources and singular potentials.

%%%%%%%%%%%%%%%%%%%%%%%%%%%%%%%%%%%%%%%%%%%%%%%%%%%%%%%%%%%%
\subsection{Potential Theory and Harmonic Analysis}
\label{subsec:potential-harmonic-analysis}
%%%%%%%%%%%%%%%%%%%%%%%%%%%%%%%%%%%%%%%%%%%%%%%%%%%%%%%%%%%%

Poisson kernels, singular integrals, maximal functions, and boundary
limits connect Potential Theory with Harmonic Analysis.

For example,

\[
u(x,t)
=
P_t*f(x)
\]

constructs a harmonic extension of boundary data.

Studying the limit

\[
u(x,t)\to f(x)
\]

as

\[
t\to0^+
\]

requires analytic tools closely related to approximate identities and
maximal operators.

%%%%%%%%%%%%%%%%%%%%%%%%%%%%%%%%%%%%%%%%%%%%%%%%%%%%%%%%%%%%
\subsection{Common Errors}
\label{subsec:potential-theory-common-errors}
%%%%%%%%%%%%%%%%%%%%%%%%%%%%%%%%%%%%%%%%%%%%%%%%%%%%%%%%%%%%

\subsubsection{Confusing Laplace and Poisson Equations}

Laplace's equation is

\[
\boxed{
\Delta u=0.
}
\]

Poisson's equation includes a source term, for example,

\[
\boxed{
-\Delta u=f.
}
\]

%%%%%%%%%%%%%%%%%%%%%%%%%%%%%%%%%%%%%%%%%%%%%%%%%%%%%%%%%%%%
\subsubsection{Ignoring Sign Conventions}

Some texts define fundamental solutions using

\[
\Delta\Phi=\delta,
\]

while others use

\[
-\Delta\Gamma=\delta.
\]

The formulas differ by a sign.

%%%%%%%%%%%%%%%%%%%%%%%%%%%%%%%%%%%%%%%%%%%%%%%%%%%%%%%%%%%%
\subsubsection{Assuming Harmonic Functions Can Have Interior Maxima}

A nonconstant harmonic function cannot attain a maximum or minimum at an
interior point of a connected domain.

%%%%%%%%%%%%%%%%%%%%%%%%%%%%%%%%%%%%%%%%%%%%%%%%%%%%%%%%%%%%
\subsubsection{Assuming Uniqueness Implies Existence}

The maximum principle proves uniqueness for many boundary-value problems,
but it does not by itself prove that a solution exists.

%%%%%%%%%%%%%%%%%%%%%%%%%%%%%%%%%%%%%%%%%%%%%%%%%%%%%%%%%%%%
\subsubsection{Forgetting Neumann Compatibility}

A Neumann problem requires an integral compatibility condition.

For

\[
\Delta u=f,
\qquad
\frac{\partial u}{\partial n}=g,
\]

one must have

\[
\boxed{
\int_\Omega f\,dx
=
\int_{\partial\Omega}g\,dS.
}
\]

%%%%%%%%%%%%%%%%%%%%%%%%%%%%%%%%%%%%%%%%%%%%%%%%%%%%%%%%%%%%
\subsubsection{Expecting Neumann Solutions to Be Unique}

If \(u\) solves a pure Neumann problem, then

\[
u+C
\]

has the same normal derivative.

Thus uniqueness is generally only up to an additive constant.

%%%%%%%%%%%%%%%%%%%%%%%%%%%%%%%%%%%%%%%%%%%%%%%%%%%%%%%%%%%%
\subsubsection{Treating Fundamental Solutions as Smooth Everywhere}

A fundamental solution is singular at its source.

The identity involving the delta distribution must be interpreted
distributionally.

%%%%%%%%%%%%%%%%%%%%%%%%%%%%%%%%%%%%%%%%%%%%%%%%%%%%%%%%%%%%
\subsubsection{Confusing Harmonic with Holomorphic}

Holomorphic functions are complex-valued functions of a complex variable.

Harmonic functions solve

\[
\Delta u=0.
\]

The real and imaginary parts of holomorphic functions are harmonic, but
harmonicity and holomorphicity are different concepts.

%%%%%%%%%%%%%%%%%%%%%%%%%%%%%%%%%%%%%%%%%%%%%%%%%%%%%%%%%%%%
\subsubsection{Ignoring Domain Geometry}

Boundary regularity, topology, and geometry can strongly influence the
existence and behavior of solutions.

%%%%%%%%%%%%%%%%%%%%%%%%%%%%%%%%%%%%%%%%%%%%%%%%%%%%%%%%%%%%
\subsubsection{Confusing Measure Zero with Zero Capacity}

A set may have measure zero but still have nonzero capacity.

Capacity captures a different notion of size.

%%%%%%%%%%%%%%%%%%%%%%%%%%%%%%%%%%%%%%%%%%%%%%%%%%%%%%%%%%%%
\subsection{Applications}
\label{subsec:potential-theory-applications}
%%%%%%%%%%%%%%%%%%%%%%%%%%%%%%%%%%%%%%%%%%%%%%%%%%%%%%%%%%%%

\begin{applicationbox}
\textbf{Electrostatics.}
Charge distributions generate electric potentials satisfying Poisson's
equation.

\medskip

\textbf{Gravitation.}
Mass distributions generate Newtonian gravitational potentials.

\medskip

\textbf{Steady-State Heat Flow.}
Temperature distributions at equilibrium are harmonic in regions without
heat sources.

\medskip

\textbf{Fluid Mechanics.}
Velocity potentials for incompressible irrotational flows are harmonic.

\medskip

\textbf{Complex Analysis.}
Real and imaginary parts of holomorphic functions are harmonic.

\medskip

\textbf{Partial Differential Equations.}
Potential theory provides model techniques for elliptic PDEs.

\medskip

\textbf{Probability.}
Harmonic functions describe expected boundary values of Brownian motion.

\medskip

\textbf{Fourier Analysis.}
Poisson kernels provide harmonic extensions and approximate identities.

\medskip

\textbf{Distribution Theory.}
Point sources and fundamental solutions are naturally interpreted as
distributions.

\medskip

\textbf{Calculus of Variations.}
Harmonic functions minimize Dirichlet energy under fixed boundary data.

\medskip

\textbf{Geometric Analysis.}
Laplace--Beltrami operators extend potential theory to curved spaces.

\medskip

\textbf{Numerical Analysis.}
Boundary-element methods use layer potentials and Green's functions to
solve boundary-value problems.
\end{applicationbox}

%%%%%%%%%%%%%%%%%%%%%%%%%%%%%%%%%%%%%%%%%%%%%%%%%%%%%%%%%%%%
\subsection{Exercises}
\label{subsec:potential-theory-exercises}
%%%%%%%%%%%%%%%%%%%%%%%%%%%%%%%%%%%%%%%%%%%%%%%%%%%%%%%%%%%%

\begin{exercise}[1]
Determine whether

\[
u(x,y)=x^2-y^2
\]

is harmonic.
\end{exercise}

\begin{exercise}[1]
Determine whether

\[
u(x,y)=x^2+y^2
\]

is harmonic.
\end{exercise}

\begin{exercise}[1]
Show that every affine function on \(\R^n\) is harmonic.
\end{exercise}

\begin{exercise}[1]
Verify that

\[
u(x,y)=xy
\]

is harmonic.
\end{exercise}

\begin{exercise}[1]
State the mean-value property for harmonic functions.
\end{exercise}

\begin{exercise}[1]
Explain why a nonconstant harmonic function cannot attain an interior
maximum.
\end{exercise}

\begin{exercise}[2]
Use the maximum principle to prove uniqueness for the Dirichlet problem.
\end{exercise}

\begin{exercise}[2]
Show that solutions of the homogeneous Neumann problem are unique only up
to constants.
\end{exercise}

\begin{exercise}[2]
Derive the compatibility condition for the Neumann problem using the
divergence theorem.
\end{exercise}

\begin{exercise}[2]
Derive Green's first identity.
\end{exercise}

\begin{exercise}[2]
Derive Green's second identity.
\end{exercise}

\begin{exercise}[2]
For a radial function

\[
u(x)=v(|x|),
\]

derive

\[
\Delta u
=
v''(r)
+
\frac{n-1}{r}v'(r).
\]
\end{exercise}

\begin{exercise}[2]
Find all radial harmonic functions on an annulus in \(\R^n\) for
\(n\geq3\).
\end{exercise}

\begin{exercise}[2]
Find all radial harmonic functions on an annulus in \(\R^2\).
\end{exercise}

\begin{exercise}[2]
Verify that

\[
|x|^{2-n}
\]

is harmonic away from the origin when \(n\geq3\).
\end{exercise}

\begin{exercise}[2]
Verify that

\[
\log|x|
\]

is harmonic away from the origin in two dimensions.
\end{exercise}

\begin{exercise}[2]
Show that

\[
u(r,\theta)=r^m\cos(m\theta)
\]

is harmonic on \(\R^2\) for every positive integer \(m\).
\end{exercise}

\begin{exercise}[2]
Find the harmonic extension to the unit disk of the boundary data

\[
g(\theta)=\sin(2\theta).
\]
\end{exercise}

\begin{exercise}[2]
Show that the Poisson kernel for the disk is positive.
\end{exercise}

\begin{exercise}[2]
Verify that the Poisson kernel has integral \(1\) around the circle.
\end{exercise}

\begin{exercise}[3]
Derive the Poisson kernel for the upper half-plane.
\end{exercise}

\begin{exercise}[3]
Derive the Poisson kernel for the unit disk using complex analysis.
\end{exercise}

\begin{exercise}[3]
Use the Dirichlet energy to derive Laplace's equation.
\end{exercise}

\begin{exercise}[3]
Prove uniqueness of the Dirichlet problem using the energy method.
\end{exercise}

\begin{exercise}[3]
Show formally that the Newtonian potential solves

\[
-\Delta u=f.
\]
\end{exercise}

\begin{exercise}[3]
Construct the Green's function for a simple one-dimensional boundary-value
problem.
\end{exercise}

\begin{exercise}[3]
Use the method of images to construct a Green's function for the upper
half-space.
\end{exercise}

\begin{exercise}[3]
Prove that the real and imaginary parts of a holomorphic function are
harmonic.
\end{exercise}

\begin{exercise}[3]
Given a harmonic function on a simply connected planar domain, derive the
equations for a harmonic conjugate.
\end{exercise}

\begin{exercise}[3]
Prove Liouville's theorem for bounded harmonic functions.
\end{exercise}

\begin{exercise}[3]
Show that a bounded isolated singularity of a harmonic function is
removable under suitable hypotheses.
\end{exercise}

\begin{exercise}[3]
Derive the Euler--Lagrange equation for the Dirichlet energy.
\end{exercise}

\begin{exercise}[3]
Explain the probabilistic interpretation of harmonic measure.
\end{exercise}

\begin{exercise}[3]
Show that the Poisson kernel defines a probability density on the
boundary.
\end{exercise}

\begin{exercise}[4]
Develop Perron's method for the Dirichlet problem.
\end{exercise}

\begin{exercise}[4]
Study the construction and properties of single-layer potentials.
\end{exercise}

\begin{exercise}[4]
Derive the jump relations for double-layer potentials on a smooth
boundary.
\end{exercise}

\begin{exercise}[4]
Study equilibrium measures for logarithmic or Newtonian energy.
\end{exercise}

\begin{exercise}[4]
Compare Lebesgue measure and capacity for thin sets.
\end{exercise}

\begin{exercise}[4]
Study the relationship between Brownian motion, harmonic measure, and the
Dirichlet problem.
\end{exercise}

\begin{exercise}[4]
Extend the mean-value property and maximum principle to suitable elliptic
operators.
\end{exercise}

\begin{exercise}[4]
Study harmonic functions on Riemannian manifolds using the
Laplace--Beltrami operator.
\end{exercise}

%%%%%%%%%%%%%%%%%%%%%%%%%%%%%%%%%%%%%%%%%%%%%%%%%%%%%%%%%%%%
\subsection{Conceptual Summary}
\label{subsec:potential-theory-summary}
%%%%%%%%%%%%%%%%%%%%%%%%%%%%%%%%%%%%%%%%%%%%%%%%%%%%%%%%%%%%

Potential theory centers on the Laplacian

\[
\boxed{
\Delta
=
\sum_{j=1}^{n}\partial_j^2.
}
\]

Harmonic functions satisfy

\[
\boxed{
\Delta u=0.
}
\]

Poisson's equation introduces sources:

\[
\boxed{
-\Delta u=f.
}
\]

Harmonic functions satisfy the mean-value property

\[
\boxed{
u(x_0)
=
\text{average of \(u\) around \(x_0\)}.
}
\]

They also satisfy maximum and minimum principles, giving strong uniqueness
results for boundary-value problems.

Fundamental solutions satisfy

\[
\boxed{
-\Delta\Gamma=\delta.
}
\]

Convolution produces potentials:

\[
\boxed{
u=\Gamma*f.
}
\]

Green's functions incorporate both

\[
\boxed{
\text{point-source behavior}
}
\]

and

\[
\boxed{
\text{boundary conditions}.
}
\]

Poisson kernels explicitly solve Dirichlet problems on highly symmetric
domains.

Harmonic functions also minimize the Dirichlet energy

\[
\boxed{
E[u]
=
\frac12
\int_\Omega|\nabla u|^2\,dx.
}
\]

Potential theory therefore connects

\[
\boxed{
\text{harmonic functions},
\quad
\text{elliptic PDEs},
\quad
\text{energy minimization},
\quad
\text{boundary values},
\quad
\text{probability}.
}
\]

%%%%%%%%%%%%%%%%%%%%%%%%%%%%%%%%%%%%%%%%%%%%%%%%%%%%%%%%%%%%
\subsection{Section Connections}
\label{subsec:potential-theory-connections}
%%%%%%%%%%%%%%%%%%%%%%%%%%%%%%%%%%%%%%%%%%%%%%%%%%%%%%%%%%%%

\chapterconnection{
Potential Theory develops the analytic structure surrounding harmonic
functions, Laplace's equation, Poisson's equation, and their fundamental
solutions. Vector Calculus supplies the gradient, divergence, Laplacian,
and Green identities. Complex Analysis connects planar harmonic functions
with holomorphic functions and harmonic conjugates. Measure Theory and
Lebesgue Integration support potential integrals and equilibrium measures.
Functional Analysis provides weak and operator-theoretic formulations of
boundary-value problems. Fourier Analysis and Harmonic Analysis provide
frequency-space and boundary-limit methods, especially through Poisson
kernels. Calculus of Variations explains harmonic functions as minimizers
of Dirichlet energy. Distribution Theory gives rigorous meaning to
fundamental solutions and point sources such as the Dirac delta. Probability
Theory interprets harmonic measure through Brownian motion. Differential
Equations and PDEs extend these ideas to general elliptic operators, while
Geometric Analysis replaces the Euclidean Laplacian with geometric
operators such as the Laplace--Beltrami operator.
}

%%%%%%%%%%%%%%%%%%%%%%%%%%%%%%%%%%%%%%%%%%%%%%%%%%%%%%%%%%%%
% SECTION SOURCE TRACKING
%%%%%%%%%%%%%%%%%%%%%%%%%%%%%%%%%%%%%%%%%%%%%%%%%%%%%%%%%%%%

%------------------------------------------------------------
% 5.15 Potential Theory
%------------------------------------------------------------
%
% Source basis:
%   - Standard classical potential theory.
%   - Standard harmonic-function theory.
%   - Standard elliptic PDE and Green-function methods.
%   - Standard variational and probabilistic interpretations.
%
% Used for:
%   - Laplace and Poisson equations;
%   - harmonic functions;
%   - mean-value properties;
%   - maximum and minimum principles;
%   - Dirichlet and Neumann problems;
%   - Green identities;
%   - fundamental solutions;
%   - radial harmonic functions;
%   - Newtonian potentials;
%   - logarithmic potentials;
%   - electrostatic and gravitational interpretations;
%   - Green's functions;
%   - representation formulas;
%   - Poisson kernels;
%   - Harnack inequalities;
%   - harmonic regularity;
%   - Liouville theorem;
%   - harmonic and holomorphic functions;
%   - subharmonic and superharmonic functions;
%   - comparison principles;
%   - barriers and Perron's method;
%   - removable singularities;
%   - layer potentials;
%   - Dirichlet energy;
%   - harmonic measure;
%   - Brownian-motion interpretation;
%   - capacity and polar sets;
%   - equilibrium measures;
%   - method of images;
%   - separation of variables;
%   - Fourier methods;
%   - spherical harmonics;
%   - potential theory on manifolds;
%   - elliptic PDE connections.
%
% Notes:
%   - Distribution Theory supplies the rigorous interpretation of
%     singular fundamental solutions and point sources.
%   - Calculus of Variations supplies the Dirichlet principle.
%   - Harmonic Analysis develops Poisson kernels, boundary convergence,
%     and singular integrals further.
%   - Probability Theory develops Brownian motion and stochastic
%     interpretations in greater depth.
%   - Differential Equations develops general elliptic boundary-value
%     problems.
%   - Geometric Analysis extends harmonic-function and elliptic theory
%     to manifolds and geometric PDEs.
%
%%%%%%%%%%%%%%%%%%%%%%%%%%%%%%%%%%%%%%%%%%%%%%%%%%%%%%%%%%%%